\section{Related Work}
\label{sec:related}

\subsection{SAR船型分类与观测、场景变化}

公开数据集为SAR船型识别提供了不同的观测基础。OpenSARShip面向Sentinel-1船舶解译建立开放数据资源\cite{huang2018opensarship}；FUSAR-Ship则将高分三号SAR影像与AIS信息匹配，覆盖多种海岸与海域场景，支持船舶检测和识别\cite{hou2020fusar}。不同数据源的传感器、图像产品、类别体系与样本组织存在差别，因此，数据集内分类表现与跨观测应用中的可靠性需要分别评价。针对数据分布差异，Zhao和Lang研究了结合双分支网络与注意力的深度子域适应方法，用于SAR船型分类\cite{zhao2022dsan}。这些研究为跨域识别提供了数据和方法基础。

本文从另一互补角度考察观测变化：局部图像证据如何参与船型判别。同船多次观测在记录身份一致的条件下呈现类内外观变化，而对指定输入分量的控制分析进一步考察分类器对这些变化的响应。这一分析着眼于类别竞争及证据利用，与特征分布适配相互补充。

\subsection{双极化融合与多尺度表征}

已有分类方法从极化互补、分层特征和融合协调等方面利用SAR船舶信息。Xie等提出的HOG-SiamShipCLSNet将HOG描述与浅层、深层卷积表征结合\cite{xie2023hogsiam}。其VV和VH分支共享骨干参数，在两个深度分别进行单极化及跨极化分组双线性池化，再与融合后的HOG特征拼接，测试时汇总多个输出。该方法联合使用双极化、显式结构描述和不同层次的学习表征，为共享特征提取与分层融合提供了直接先例。

DPIG-Net采用VV主导的三分支组织\cite{shao2023dpig}：VV幅度作为主输入，VH幅度与跨极化乘积幅度作为辅助输入；交叉注意力利用VV特征引导辅助分支，经拼接后由膨胀残差稠密结构融合，并通过VV全局残差保留主分支信息。其设计同时考虑极化交互、上下文聚合，以及融合过程中对主分支信息的保留。

CPINet进一步结合多尺度深层特征融合、混合阶统计注意力与因子化双线性编码\cite{he2024cpinet}。MO-SE使用一路极化的一阶、二阶统计生成权重，调制另一路特征；FBC既聚合跨极化关系，也处理单极化表征，并通过自适应损失权重协调四个分类输出。其伪孪生特征提取与共享FC嵌入层的组织方式，也体现了极化专属处理与共享表征之间的结构选择。

上述方法分别在不同类别配置及数据划分下验证了融合设计：HOG-SiamShipCLSNet使用OpenSARShip 2.0四类任务，DPIG-Net报告三类和六类实验，CPINet报告三类和五类实验\cite{xie2023hogsiam,shao2023dpig,he2024cpinet}。其结果为相应设定中的极化融合提供了依据，跨方法比较则需要统一类别体系、预处理和评价协议。

在已有融合研究的基础上，本文进一步分析这些信息在观测变化下的具体判别用途：某一极化、区域和尺度的响应何时支持正确类别，何时加强混淆。区域用于区分船体及邻域的空间来源，尺度用于刻画不同层次的图像或特征变化。保留这些身份，使局部响应的变化能够与具体类别竞争联系起来，并为后续证据组织方法提供可检验的设计依据。

\subsection{受控扰动、解释与训练约束}

通过输入修改分析模型依赖已有明确方法基础。Li等针对MSTAR上的SAR自动目标识别，利用Shapley值量化目标、阴影和杂波区域的贡献，分析区域间交互，并通过随机调整信杂比研究背景依赖\cite{li2023noncausality}。这项工作已将数据选择偏差、区域信息和模型行为联系起来，而非仅展示注意力图。但其地面目标结论不能直接证明船舶周边响应都属于无用背景，也不能确定某一次船型误判的物理原因。

Meaningful Perturbation通过优化紧凑掩码，以模糊、替换等显式操作压低或保留分类响应，并用正则化限制扰动伪影\cite{fong2017meaningful}。这为研究固定模型的输入依赖提供了依据。ROAR则将重要性排序对应的特征从训练集和测试集中同时移除，再重新训练，以减少仅在测试时删除信息引入的分布变化对解释评估的干扰\cite{hooker2019roar}。二者回答的问题不同：固定模型对某次修改敏感，不等于删除该信息后重新训练必然变差，也不等于据此增加一个模块就能提高完整模型表现。

Yin等的Fourier Perspective使用范数受控的傅里叶基扰动考察不同频率位置的敏感性，并分析不同图像扰动之间的鲁棒性权衡\cite{yin2019fourier}。它提示应区分扰动强度与频率位置，而不能预设高频均为噪声。对于SAR灰度图像，空间尺度或傅里叶成分也不能直接命名为某种雷达散射、运动或环境机制。本文据此保留干预的极化、区域与尺度身份，同时检查正确类和实际竞争类的分数变化；救回可能来自混淆类别受到更强抑制，并不总是正确类响应增强。

解释信息也可以进入训练目标。Ross等的Right for the Right Reasons通过惩罚被标注为无关位置的输入梯度，将解释约束加入分类训练\cite{ross2017right}。这一工作建立了解释与可微学习之间的联系。对于局部证据作用随样本和联合读出变化的船型识别，值得进一步研究的是如何从配对扰动中提取可靠的判别约束：既利用能够减轻类别混淆的变化，也保护原本支持正确判断的证据。此类约束的价值应通过与普通增强及匹配训练对照的比较检验，而不能仅由扰动敏感性推定。

